% Default EPTCS-oriented preamble for md2tex Workshop.
% This preamble intentionally omits title/author/event metadata.

\PassOptionsToPackage{hidelinks}{hyperref}
\documentclass{eptcs}

\usepackage{amsmath}
\usepackage{amssymb}
\usepackage{mathtools}
\usepackage{amsthm}
\usepackage{braket}

\usepackage{graphicx}
\usepackage{booktabs}
\usepackage{xcolor}
\usepackage{float}
\usepackage{aliascnt}

\IfFileExists{algorithm.sty}{%
  \usepackage{algorithm}
}{}
\IfFileExists{algpseudocode.sty}{%
  \usepackage{algpseudocode}
  \algrenewcommand\algorithmicrequire{Input:}
  \algrenewcommand\algorithmicensure{Output:}
  \algnewcommand\algorithmicinput{\textbf{Input:}}
  \algnewcommand\Input{\item[\algorithmicinput]}
  \algnewcommand\algorithmicoutput{\textbf{Output:}}
  \algnewcommand\Output{\item[\algorithmicoutput]}
}{}

\IfFileExists{bbold.sty}{%
  \usepackage{bbold}
  \newcommand{\indicatorone}{\mathbb{1}}
}{%
  \newcommand{\indicatorone}{\mathbf{1}}
}

\newtheorem{theorem}{Theorem}[section]

\newaliascnt{lemma}{theorem}
\newtheorem{lemma}[lemma]{Lemma}
\aliascntresetthe{lemma}

\newaliascnt{proposition}{theorem}
\newtheorem{proposition}[proposition]{Proposition}
\aliascntresetthe{proposition}

\newaliascnt{corollary}{theorem}
\newtheorem{corollary}[corollary]{Corollary}
\aliascntresetthe{corollary}

\theoremstyle{definition}
\newaliascnt{definition}{theorem}
\newtheorem{definition}[definition]{Definition}
\aliascntresetthe{definition}

\newaliascnt{example}{theorem}
\newtheorem{example}[example]{Example}
\aliascntresetthe{example}

\theoremstyle{remark}
\newaliascnt{remark}{theorem}
\newtheorem{remark}[remark]{Remark}
\aliascntresetthe{remark}

\newaliascnt{property}{theorem}
\newtheorem{property}[property]{Property}
\aliascntresetthe{property}

\IfFileExists{cleveref.sty}{%
  \usepackage[nameinlink,capitalize,noabbrev]{cleveref}
  \crefname{lemma}{Lemma}{Lemmas}
  \Crefname{lemma}{Lemma}{Lemmas}
  \crefname{proposition}{Proposition}{Propositions}
  \Crefname{proposition}{Proposition}{Propositions}
  \crefname{corollary}{Corollary}{Corollaries}
  \Crefname{corollary}{Corollary}{Corollaries}
  \crefname{definition}{Definition}{Definitions}
  \Crefname{definition}{Definition}{Definitions}
  \crefname{example}{Example}{Examples}
  \Crefname{example}{Example}{Examples}
  \crefname{remark}{Remark}{Remarks}
  \Crefname{remark}{Remark}{Remarks}
  \crefname{property}{Property}{Properties}
  \Crefname{property}{Property}{Properties}
  \crefformat{equation}{(##2##1##3)}
  \Crefformat{equation}{(##2##1##3)}
  \crefrangeformat{equation}{(##3##1##4)--(##5##2##6)}
  \Crefrangeformat{equation}{(##3##1##4)--(##5##2##6)}
  \crefmultiformat{equation}{(##2##1##3)}{, (##2##1##3)}{, (##2##1##3)}{, (##2##1##3)}
  \Crefmultiformat{equation}{(##2##1##3)}{, (##2##1##3)}{, (##2##1##3)}{, (##2##1##3)}
}{%
  \providecommand{\cref}[1]{\ref{#1}}
  \providecommand{\Cref}[1]{\ref{#1}}
}

\newcommand{\N}{\mathbb{N}}
\newcommand{\Z}{\mathbb{Z}}
\newcommand{\Q}{\mathbb{Q}}
\newcommand{\R}{\mathbb{R}}
\newcommand{\C}{\mathbb{C}}

\providecommand{\set}[1]{\left\{#1\right\}}
\newcommand{\abs}[1]{\left|#1\right|}
\newcommand{\norm}[1]{\left\lVert#1\right\rVert}

\providecommand{\keywords}[1]{\par\noindent\textbf{Keywords:} #1\par}
\newcommand{\ind}[1]{\indicatorone_{\{#1\}}}
\newcommand{\1}{\indicatorone}

\bibliographystyle{eptcs}
